% from ega2-1-fr.tex
\oldpage[II]{10}
\begin{proposition}[1.4.3]
\label{II.1.4.3-fr}
Pour tout couple $(\mathcal{B},\mathcal{M})$ formé d'une
$\mathcal{O}_S$-Algèbre quasi-cohérente $\mathcal{B}$ et d'un
$\mathcal{B}$-Module $\mathcal{M}$ quasi-cohérent (en tant que
$\mathcal{O}_S$-Module ou en tant que $\mathcal{B}$-Module, ce qui revient au
même (\hyperref[I.9.6.1-fr]{I, 9.6.1})), il existe un couple
$(X,\mathcal{F})$ formé d'un préschéma $X$ affine au-dessus de $S$ et d'un
$\mathcal{O}_X$-Module quasi-cohérent $\mathcal{F}$, tels que
$\mathcal{A}(X)=\mathcal{B}$ et $\mathcal{A}(\mathcal{F})=\mathcal{M}$~; en
outre ce couple est déterminé à un isomorphisme unique près.
\end{proposition}

L'unicité résulte de (\hyperref[II.1.4.1-fr]{1.4.1}) et
(\hyperref[II.1.4.2-fr]{1.4.2})~; l'existence se démontre comme dans
(\hyperref[II.1.3.1-fr]{1.3.1}), et nous laissons encore les détails au
lecteur. On note $\widetilde{\mathcal{M}}$ le $\mathcal{O}_X$-Module
$\mathcal{F}$, et on dit qu'il est \emph{associé} au $\mathcal{B}$-Module
quasi-cohérent $\mathcal{M}$~; pour tout ouvert affine $U\subset S$,
$\widetilde{\mathcal{M}}|p^{-1}(U)$ (où $p$ est le morphisme structural
$X\to S$) est canoniquement isomorphe à
$(\Gamma(U,\mathcal{M}))^{\sim}$.

\begin{corollary}[1.4.4]
\label{II.1.4.4-fr}
Dans la catégorie des $\mathcal{B}$-Modules quasi-cohérents,
$\widetilde{\mathcal{M}}$ est un foncteur covariant additif exact en
$\mathcal{M}$, qui commute aux limites inductives et aux sommes directes.
\end{corollary}

On est en effet aussitôt ramené au cas où $S$ est affine, et le corollaire
résulte alors de (\hyperref[I.1.3.5-fr]{I, 1.3.5}),
(\hyperref[I.1.3.9-fr]{I, 1.3.9}) et
(\hyperref[I.1.3.11-fr]{I, 1.3.11}).

\begin{corollary}[1.4.5]
\label{II.1.4.5-fr}
Sous les hypothèses de (\hyperref[II.1.4.3-fr]{1.4.3}), pour que
$\widetilde{\mathcal{M}}$ soit un $\mathcal{O}_X$-Module de type fini, il faut
et il suffit que $\mathcal{M}$ soit un $\mathcal{B}$-Module de type fini.
\end{corollary}

On est aussitôt ramené au cas où $S=\operatorname{Spec}(A)$ est un schéma
affine. Alors $\mathcal{B}=\widetilde{B}$, où $B$ est une $A$-algèbre de type
fini (\hyperref[I.9.6.2-fr]{I, 9.6.2}), et
$\mathcal{M}=\widetilde{M}$, où $M$ est un $B$-module
(\hyperref[I.1.3.13-fr]{I, 1.3.13})~; \emph{sur le préschéma $X$},
$\mathcal{O}_X$ est associé à l'anneau $B$ et
$\widetilde{\mathcal{M}}$ au $B$-module $M$~; pour que
$\widetilde{\mathcal{M}}$ soit de type fini, il faut et il suffit donc que
$M$ soit de type fini (\hyperref[I.1.3.13-fr]{I, 1.3.13}), d'où notre
assertion.

\begin{proposition}[1.4.6]
\label{II.1.4.6-fr}
Soient $Y$ un préschéma affine au-dessus de $S$, $X$, $X'$ deux préschémas
affines au-dessus de $Y$ (donc aussi au-dessus de $S$
(\hyperref[II.1.3.5-fr]{1.3.5})). Soient
$\mathcal{B}=\mathcal{A}(Y)$, $\mathcal{A}=\mathcal{A}(X)$,
$\mathcal{A}'=\mathcal{A}(X')$. Alors $X\times_Y X'$ est affine au-dessus de
$Y$ (donc aussi au-dessus de $S$) et $\mathcal{A}(X\times_Y X')$
s'identifie canoniquement à $\mathcal{A}\otimes_{\mathcal{B}}\mathcal{A}'$.
\end{proposition}

En effet (\hyperref[I.9.6.1-fr]{I, 9.6.1})
$\mathcal{A}\otimes_{\mathcal{B}}\mathcal{A}'$ est une
$\mathcal{B}$-Algèbre quasi-cohérente, donc aussi une
$\mathcal{O}_S$-Algèbre quasi-cohérente
(\hyperref[I.9.6.1-fr]{I, 9.6.1})~; soit $Z$ le spectre de
$\mathcal{A}\otimes_{\mathcal{B}}\mathcal{A}'$
(\hyperref[II.1.3.1-fr]{1.3.1}). Les $\mathcal{B}$-homomorphismes canoniques
$\mathcal{A}\to\mathcal{A}\otimes_{\mathcal{B}}\mathcal{A}'$ et
$\mathcal{A}'\to\mathcal{A}\otimes_{\mathcal{B}}\mathcal{A}'$ correspondent
(\hyperref[II.1.2.7-fr]{1.2.7}) à des $Y$-morphismes $p:Z\to X$ et
$p':Z\to X'$. Pour voir que le triplet $(Z,p,p')$ est un produit
$X\times_Y X'$, on peut se borner au cas où $S$ est un schéma affine
d'anneau $C$ (\hyperref[I.3.2.6.4-fr]{I, 3.2.6.4}). Mais alors $Y$, $X$,
$X'$ sont des schémas affines (\hyperref[II.1.3.4-fr]{1.3.4}) dont les
anneaux $B$, $A$, $A'$ sont des $C$-algèbres telles que
$\mathcal{B}=\widetilde{B}$, $\mathcal{A}=\widetilde{A}$,
$\mathcal{A}'=\widetilde{A'}$. On sait alors
(\hyperref[I.1.3.13-fr]{I, 1.3.13}) que
$\mathcal{A}\otimes_{\mathcal{B}}\mathcal{A}'$ s'identifie canoniquement à
la $\mathcal{O}_S$-Algèbre $\widetilde{A\otimes_B A'}$, donc l'anneau
$A(Z)$ s'identifie à $A\otimes_B A'$ et les morphismes $p$, $p'$ correspondent
aux homomorphismes canoniques $A\to A\otimes_B A'$,
$A'\to A\otimes_B A'$. La proposition résulte alors de
(\hyperref[I.3.2.2-fr]{I, 3.2.2}).

\begin{corollary}[1.4.7]
\label{II.1.4.7-fr}
Soit $\mathcal{F}$ (resp. $\mathcal{F}'$) un $\mathcal{O}_X$-Module (resp. un
$\mathcal{O}_{X'}$-Module) quasi-cohérent~; alors
$\mathcal{A}(\mathcal{F}\otimes_Y\mathcal{F}')$ s'identifie canoniquement à
$\mathcal{A}(\mathcal{F})\otimes_{\mathcal{A}(Y)}
\mathcal{A}(\mathcal{F}')$.
\end{corollary}

On sait que $\mathcal{F}\otimes_Y\mathcal{F}'$ est quasi-cohérent sur
$X\times_Y X'$ (\hyperref[I.9.1.2-fr]{I, 9.1.2}). Soient $g:Y\to S$,
$f:X\to Y$, $f':X'\to Y$ les morphismes structuraux, de sorte que le
morphisme structural
